\chapter{Algoritma Komputasi Frekuensi Eigen Schwarzschild Anti-de Sitter}
\label{app:komputasi-sads}

\lstset{
    language=Python,
    basicstyle=\ttfamily\small,
    breaklines=true,
    breakatwhitespace=false,
    breakautoindent=true,
    postbreak=\mbox{\textcolor{gray}{$\hookrightarrow$}\space},
    columns=fullflexible,
    keepspaces=true,
    showstringspaces=false,
    tabsize=4,
    frame=none
}

\begin{lstlisting}[language=Python]
import numpy as np
import matplotlib.pyplot as plt

def _stu(r_plus, l, omega, m2):
    xp = 1.0/r_plus; c = l*(l+1)
    s = np.array([xp*(xp**2+3), 4*xp**2+9, 6*xp+10/xp, 4+5/xp**2,
                  (xp**2+1)/xp**3], dtype=complex)
    t = np.array([xp*(xp**2+3)-2j*omega*xp**2, 6*xp**2+12-4j*omega*xp,
                  12*xp+18/xp-2j*omega, 10+12/xp**2, 3/xp+3/xp**3], dtype=complex)
    u = np.array([0, (c+1)*xp**2+3+m2, 2*c*xp+3*xp+3/xp, c+3+3/xp**2,
                  (xp**2+1)/xp**3], dtype=complex)
    return s, t, u

def FA(omega, r_plus, l, m2, N=800):
    xp = 1.0/r_plus
    kappa = 0.5*(2/xp + (r_plus**3+r_plus)*xp**2)
    s, t, u = _stu(r_plus, l, omega, m2)
    b = np.zeros(N+1, dtype=complex); b[0] = 1.0
    for nn in range(1, N+1):
        Pn = 2*xp**2*nn*(nn*kappa - 1j*omega)
        acc = 0j
        for k in range(max(0, nn-4), nn):
            m = nn-k
            acc += (k*(k-1)*s[m] + k*t[m] + u[m]) * ((-xp)**m) * b[k]
        b[nn] = -acc/Pn
    return np.sum(b)

def findA(r_plus, l, m2, guess, N=800, tol=1e-10, it=200):
    f = lambda w: FA(w, r_plus, l, m2, N)
    z, h = complex(guess), 1e-7
    for _ in range(it):
        fz = f(z)
        if abs(fz) < tol:
            return z, abs(fz)
        fp = (f(z+h) - f(z-h))/(2*h)
        if abs(fp) < 1e-300:
            break
        zn = z - fz/fp
        if abs(zn - z) < tol:
            return zn, abs(f(zn))
        z = zn
    return z, abs(f(z))

def overtones0(r_plus, l, nmax=3, N=800):
    g0 = r_plus*(1.85-2.66j)
    found = [findA(r_plus, l, 0.0, g0, N)[0]]
    found.append(findA(r_plus, l, 0.0, found[0]+r_plus*(1.31-2.26j), N)[0])
    for _ in range(2, nmax):
        found.append(findA(r_plus, l, 0.0, found[-1]+(found[-1]-found[-2]), N)[0])
    return found

def to_mass(r_plus, l, om0, m2_target=1.0, dstep=0.25, N=800):
    g = om0
    for x in np.linspace(0, m2_target, int(np.ceil(m2_target/dstep))+1)[1:]:
        om, _ = findA(r_plus, l, x, g, N); g = om
    return om

LS, NS = [0, 1, 2], [0, 1, 2]
COLORS = {0: "tab:blue", 1: "tab:red", 2: "tab:green"}

# Gunakan grid komputasi rapat agar pencarian akar stabil saat turun,
# tetapi hanya simpan nilai r+ diskret yang diminta.
CALC_STEPS = [100.0, 50.0, 10.0, 5.0, 2.0, 1.0, 0.8, 0.6, 0.4]
TARGET_RS = [0.4, 0.8, 1.0, 10.0, 50.0, 100.0]

# ---- hitung mu=0 (continuation) lalu mu=1 ----
DATA = {}
for l in LS:
    ov = overtones0(CALC_STEPS[0], l)
    seeds = {n: ov[n] for n in NS}
    rprev = CALC_STEPS[0]

    for rp in CALC_STEPS:
        for n in NS:
            # Scaling tebakan awal: proporsional r+ jika rp >= 1, konstan jika rp < 1
            if rp >= 1.0:
                guess = seeds[n] * (rp / rprev)
            else:
                guess = seeds[n]

            om, _ = findA(rp, l, 0.0, guess)
            seeds[n] = om

            if rp in TARGET_RS:
                DATA[(l, n, rp, 0.0)] = om
                DATA[(l, n, rp, 1.0)] = to_mass(rp, l, om)
        rprev = rp

RS = np.array(TARGET_RS)

# ---- 6 gambar terpisah ----
for n in NS:
    for part, ylab in [("R", r"$\omega_R$"),
                       ("I", r"$\omega_I$")]:
        plt.figure(figsize=(7.8, 5.6))
        for l in LS:
            y0 = [(DATA[(l, n, rp, 0.0)].real if part == "R"
                   else -DATA[(l, n, rp, 0.0)].imag) for rp in RS]
            y1 = [(DATA[(l, n, rp, 1.0)].real if part == "R"
                   else -DATA[(l, n, rp, 1.0)].imag) for rp in RS]

            plt.plot(RS, y0, "-o", color=COLORS[l], ms=5, label=fr"$\ell={l},\ \mu=0$")
            plt.plot(RS, y1, "--s", color=COLORS[l], ms=5, mfc="white",
                     label=fr"$\ell={l},\ \mu=1$")

        # Skala linier, label disesuaikan
        plt.xlabel(r"$r_+$")
        plt.ylabel(ylab)
        plt.grid(alpha=0.3)
        plt.legend(fontsize=8)
        plt.tight_layout()
        plt.savefig(f"grafik_sads_n{n}_w{part}.png", dpi=140, bbox_inches="tight")
        plt.close()
print("Tersimpan 6 gambar SAdS: grafik_sads_n{0,1,2}_{wR,wI}.png")

# ---- tabel ----
print("\nTABEL SAdS - omega = omega_R - i*omega_I")
for n in NS:
    print(f"\n  overtone n = {n}")
    print(f"  {'l':>2} {'r+':>5} | {'wR(mu0)':>9} {'wI(mu0)':>9} | {'wR(mu1)':>9} {'wI(mu1)':>9}")
    print("  "+"-"*52)
    for l in LS:
        for rp in RS:
            a = DATA[(l, n, rp, 0.0)]; b = DATA[(l, n, rp, 1.0)]
            print(f"  {l:>2} {rp:>5.1f} | {a.real:>9.3f} {-a.imag:>9.3f} | {b.real:>9.3f} {-b.imag:>9.3f}")
\end{lstlisting}